\subsection*{VI.21.P16: Pointwise residue characteristic versus scheme-theoretic characteristic}
\begin{problem}\label{prob:vi21-p16}
Assume every point $x\in X$ has residue field of characteristic $p$.  Show that the underlying
point map $|X|\to|\Spec\mathbb Z|$ has image contained in $\{(p)\}$.  Show by example that this
does not imply that the structure morphism factors through $\Spec\mathbb F_p$.  State the correct
factorization criterion.
\par\textbf{Solution.} For each $x$, the image point is
$\ker(\mathbb Z\to\kappa(x))=(p)$, so the point-set image is the singleton $\{(p)\}$.
However take $X=\Spec(\mathbb Z/p^2\mathbb Z)$.  Its unique prime is $(p)/(p^2)$ and its residue
field is $\mathbb F_p$, but $p\cdot1\neq0$ in $\mathbb Z/p^2\mathbb Z$.  Hence there is no unital
ring map $\mathbb F_p\to\mathbb Z/p^2\mathbb Z$ compatible with $\mathbb Z$, so the structure
morphism does not factor through $\Spec\mathbb F_p$.  The correct criterion is exactly
$p\cdot1=0$ in $\Gamma(X,\OO_X)$, equivalently the integer section $p\cdot1$ is zero in every
stalk.
\end{problem}
\paragraph{Provenance.} Canonical VI/21 derivative/application; not an additional legacy migration.
